% PDF demonstrativo público. Todos os nomes, números e fatos são fictícios.
% A opção permitir-placeholders é exclusiva desta demonstração e dos testes.
\documentclass[permitir-placeholders]{sindicancia}

\dataPublicacaoPortariaISO{20260501}
\dataRecebimentoPortariaISO{20260502}
\ritoSindicancia{sumario}
\naturezaSindicancia{processual}
\partePrincipalSindicancia{sindicado}
\classificacaoSindicancia{ostensiva}

\dia{2}
\mes{maio}
\ano{2026}
\cidade{Salvador/BA}

\om{1º Batalhão de Demonstração}
\omSuperior{Comando de Projetos Demonstrativos}
\omComplemento{Unidade fictícia para documentação}
\omNomeHistorico{Unidade-Modelo}

\nup{00001.000001/2026-11}
\portaria{Portaria nº 12, de 1º de maio de 2026, do Comandante do 1º Batalhão de Demonstração}
\objeto{apurar, em cenário inteiramente fictício, divergência no registro de entrega de material de expediente}

\autoridadeNome{Carlos Exemplo}
\autoridadePosto{Coronel}
\autoridadeFuncao{Comandante do 1º Batalhão de Demonstração}

\sindicanteNome{Ana Modelo}
\sindicantePosto{Capitão}
\sindicanteOM{1º Batalhão de Demonstração}

\sindicadoNome{João Fictício}
\sindicadoPosto{2º Sargento}
\sindicadoSU{Seção Administrativa}
\sindicadoOM{1º Batalhão de Demonstração}
\sindicadoIdentidade{000000000 MD/EB}
\sindicadoCPF{000.000.000-00}
\sindicadoNascimento{1º de janeiro de 1990}
\sindicadoNaturalidade{Salvador/BA}
\sindicadoEstadoCivil{casado}
\sindicadoFiliacao{Pedro Fictício e Maria Fictícia}
\sindicadoResidencia{Rua do Exemplo, nº 100, Salvador/BA}
\sindicadoCelular{(00) 00000-0000}

\escrivaoNome{Marcos Demonstração}
\escrivaoPosto{1º Tenente}
\prazo{3 (três) dias úteis}

\AddToHook{shipout/background}{%
  \begin{tikzpicture}[remember picture,overlay]
    \node[rotate=45,scale=4.2,text opacity=0.08] at (current page.center)
      {DOCUMENTO FICTÍCIO -- SEM VALIDADE};
  \end{tikzpicture}%
}

\begin{document}

\pdfbookmark[0]{Guia demonstrativo dos anexos}{guia-demonstrativo}
\thispagestyle{empty}
\cabecalhoInstitucional
\vspace{1.5em}
\begin{center}
  {\Large\textbf{DEMONSTRAÇÃO COMPLETA DOS ANEXOS, DO A AO AA}}\\[0.5em]
  {\large EB10-IG-09.001 - versão 2.0.0 do projeto}
\end{center}
\vspace{1em}

Este PDF é a porta de entrada visual do projeto. Ele reúne um exemplo fictício de
cada anexo implementado, preservando a ordem normativa. Os anexos A, B, Z e AA
são mostrados apenas como referência da Seção de Apoio Jurídico e não são
elaborados pelo sindicante.

\vspace{1em}
\noindent
\begin{minipage}[t]{0.48\textwidth}
\textbf{Anexos A a N}
\begin{itemize}
  \item[A] Portaria de instauração
  \item[B] Portaria por denúncia anônima
  \item[C] Capa dos autos
  \item[D] Termo de abertura
  \item[E] Juntada inicial
  \item[F] Designação de escrivão
  \item[G] Compromisso de escrivão
  \item[H] Despacho
  \item[I] Notificação prévia
  \item[J] Notificação de novo sindicado
  \item[K] Diligências complementares
  \item[L] Juntada posterior
  \item[M] Comparecimento de sindicado
  \item[N] Comparecimento de testemunha
\end{itemize}
\end{minipage}\hfill
\begin{minipage}[t]{0.48\textwidth}
\textbf{Anexos O a AA}
\begin{itemize}
  \item[O] Carta precatória
  \item[P] Inquirição de testemunha
  \item[Q] Inquirição de sindicado
  \item[R] Substituição de sindicante
  \item[S] Acareação
  \item[T] Encerramento da instrução
  \item[U] Vista para alegações finais
  \item[V] Certidão
  \item[W] Relatório
  \item[X] Termo de encerramento
  \item[Y] DIEx de remessa
  \item[Z] Solução
  \item[AA] Modificação do rito
\end{itemize}
\end{minipage}

\vfill
\begin{center}
\textbf{Material exclusivamente demonstrativo - sem validade administrativa.}
\end{center}
\clearpage

% ANEXO A - referência da Seção de Apoio Jurídico
% DADOS - ANEXO A
\dataDocumento{1}{maio}{2026}{CIDADE/UF}
\chaveCronologica{20260501}
\modoDocumento{padrao}
\portariaNumero{001/2026}
\portariaAnexoOrigem{DIEx nº 001, de 30 de abril de 2026}
\portariaSinteseFatos{fato de interesse da Administração Militar descrito no documento anexo}
\portariaPrazoDias{30}
% Opcional, apenas quando a complexidade recomendar:
% \portariaEscrivao{POSTO/GRADUAÇÃO NOME DO ESCRIVÃO}
\portariaPublicacaoBI{BI nº 080, de 1º de maio de 2026}

\GerarDocumento{portaria-instauracao}

% ANEXO B - referência da Seção de Apoio Jurídico
% DADOS - ANEXO B
% Referência da Seção de Apoio Jurídico; não é elaborado pelo sindicante.
\dataDocumento{1}{maio}{2026}{CIDADE/UF}
\chaveCronologica{20260501}
\modoDocumento{padrao}
\portariaNumero{002/2026}
\portariaAnexoOrigem{Relatório de verificação sumária nº 001/2026}
\portariaSinteseFatos{irregularidade constatada em verificação preliminar}
\portariaPrazoDias{30}
\portariaPublicacaoBI{BI nº 080, de 1º de maio de 2026}
\portariaAnonimaDocumentoVerificacao{Relatório de verificação sumária nº 001/2026}
\portariaAnonimaEnvolvido{}
\portariaAnonimaOM{}
\portariaAnonimaIrregularidade{fatos verificados de forma independente da denúncia anônima}
\portariaAnonimaInteressePublico{o interesse do Exército no adequado esclarecimento dos fatos}
\confirmarDenunciaAnonimaNaoJuntada

\GerarDocumento{portaria-denuncia-anonima}

% ANEXO C
% ============================================================
% DADOS - ANEXO C - CAPA DOS AUTOS
% Normalmente estes campos ficam em dados/dados-sindicancia.tex.
% ============================================================

\nup{64108.000000/2026-00}
\objeto{apurar os fatos constantes da documentação inicial da sindicância.}

\sindicanteNome{RAFAEL DOS SANTOS DANTAS}
\sindicantePosto{1º Ten}

% Preencha apenas se houver sindicado definido.
\sindicadoNome{NOME DO SINDICADO}
\sindicadoPosto{POSTO/GRADUAÇÃO}

\GerarDocumento{capa}

% ANEXO D
% ============================================================
% DADOS - ANEXO D - TERMO DE ABERTURA
% Arquivo: dados/dados-anexo-d-termo-abertura.tex
%
% Este anexo usa apenas dados gerais definidos em:
%   dados/dados-sindicancia.tex
%
% Não redefina aqui macros como:
%   \dia
%   \mes
%   \ano
%   \cidade
%   \om
%   \portaria
%   \sindicanteNome
%   \sindicantePosto
%
% O arquivo é mantido para preservar a organização por anexo
% e permitir comentários específicos do termo de abertura, se necessário.
% ============================================================

\GerarDocumento{termo-abertura}

% ANEXO E
% ============================================================
% DADOS - ANEXOS E/L - JUNTADA 1
% Salvar como dados/dados-juntada1.tex para uso real.
% ============================================================

\juntadaNumero{1}

% Padrão automático: documentos/docjuntada1.pdf
% Para mudar o arquivo:
% \juntadaArquivo{documentos/outro-arquivo.pdf}

\juntadaItens{%
  \item Portaria nº 000, de 00 de mês de 2026, do Sr. Comandante do 72º Batalhão de Infantaria de Caatinga.
  \item DIEx nº 000, de 00 de mês de 2026, do Sr. NOME/POSTO/FUNÇÃO.
}

\GerarDocumento{juntada}

% ANEXO F
% ============================================================
% DADOS: ANEXO F - DESIGNAÇÃO DE ESCRIVÃO
% Arquivo: dados/dados-anexo-f-designacao-escrivao.tex
% ============================================================

\designacaoEscrivaoNumero{1}

% Se este campo ficar vazio, o modelo usa \portaria{} de dados-sindicancia.tex.
% \designacaoEscrivaoPortaria{Portaria nº ..., de ...}

% Se ficar vazio, o modelo usa \localData.
% \designacaoEscrivaoLocalData{Petrolina/PE, 31 de maio de 2026}

\escrivaoNome{NOME COMPLETO DO ESCRIVÃO}
\escrivaoPosto{POSTO/GRADUAÇÃO}

\GerarDocumento{designacao-escrivao}

% ANEXO G
% ============================================================
% DADOS: ANEXO G - COMPROMISSO DE ESCRIVÃO
% Arquivo: dados/dados-anexo-g-compromisso-escrivao.tex
% ============================================================

% Se ficar vazio, o modelo usa \localData.
% \compromissoEscrivaoLocalData{Petrolina/PE, 31 de maio de 2026}

\escrivaoNome{NOME COMPLETO DO ESCRIVÃO}
\escrivaoPosto{POSTO/GRADUAÇÃO}

\GerarDocumento{compromisso-escrivao}

% ANEXO H
% ============================================================
% DADOS - ANEXO H - DESPACHO 1
% Salvar como dados/dados-despacho1.tex para uso real.
% ============================================================

\novoDespacho

\despachoNumero{1}

\despachoTexto{%
Oficiar ao Sr. Comandante do(a) ORGANIZAÇÃO MILITAR, solicitando a remessa de cópia de DOCUMENTO/INFORMAÇÃO necessária à presente sindicância.
}

% Se quiser data própria:
% \despachoLocalData{Petrolina/PE, 31 de maio de 2026}

\GerarDocumento{despacho}

% ANEXO I
% ============================================================
% DADOS - ANEXO I - NOTIFICAÇÃO PRÉVIA 1
% Salvar como dados/dados-notificacao-previa1.tex para uso real.
% ============================================================

% ------------------------------------------------------------
% Data própria do documento
% Deixe em branco para usar a data/local do termo de abertura.
% Formato: \dataDocumento{dia}{mês}{ano}{cidade/UF}
% ------------------------------------------------------------
\dataDocumento{5}{maio}{2026}{CIDADE/UF}
\chaveCronologica{20260505}

\novaNotificacaoPrevia

\notificacaoNumero{001}

\notificacaoAnexos{cópia da \getPortaria; cópia dos documentos que deram origem à instauração.}
\notificacaoLocalVista{Sala do Sindicante, no quartel do \getOm}
\notificacaoCanal{direta}
\notificacaoAntecedenciaDiasUteis{3}
\notificacaoFatosEspecificos{possível irregularidade descrita no objeto da sindicância, nas circunstâncias de tempo, lugar e modo constantes dos documentos anexos}
\notificacaoFormaCiencia{recibo pessoal datado e assinado ao final deste documento}

\localInquiricao{Sala do Sindicante, no quartel do \getOm}
\dataInquiricao{05 de junho de 2026}
\horaInquiricao{08h00}

% Opcional: quadro de testemunhas/ofendido.
% \notificacaoTestemunhas{%
% 1 & 2º Sgt FULANO DE TAL & Sala do Sindicante & 05 JUN 26, às 09h00 \\ \hline
% 2 & Cb CICLANO DE TAL & Sala do Sindicante & 05 JUN 26, às 10h00 \\ \hline
% }

\GerarDocumento{notificacao-previa}

% ANEXO J
% ============================================================
% DADOS: ANEXO J - NOTIFICAÇÃO DE ENVOLVIDO QUE PASSA À CONDIÇÃO DE SINDICADO
% Arquivo: dados/dados-anexo-j-notificacao-novo-sindicado.tex
% ============================================================

% ------------------------------------------------------------
% Data própria do documento
% Deixe em branco para usar a data/local do termo de abertura.
% Formato: \dataDocumento{dia}{mês}{ano}{cidade/UF}
% ------------------------------------------------------------
\dataDocumento{6}{maio}{2026}{CIDADE/UF}
\chaveCronologica{20260506}

\novaNotificacaoNovoSindicado

\notificacaoNovoSindicadoNumero{001}
\notificacaoNovoSindicadoOrigem{Sindicante}

% Se ficar vazio, o modelo usa \localData.
% \notificacaoNovoSindicadoLocalData{CIDADE/UF, 31 de maio de 2026}

\notificacaoNovoSindicadoDestinatario{\getSindicadoPosto\ \getSindicadoNome\ -- \getSindicadoSU/\getSindicadoOM}
\notificacaoNovoSindicadoAnexos{cópia da \getPortaria; cópia dos documentos que deram origem à instauração}
\notificacaoNovoSindicadoLocalVista{Sala do Sindicante, no quartel do \getOm}
\notificacaoCanal{direta}
\notificacaoAntecedenciaDiasUteis{3}
\notificacaoFatosEspecificos{elementos apurados que indicam a possível participação do destinatário nos fatos delimitados pela portaria instauradora}
\notificacaoFormaCiencia{recibo pessoal datado e assinado ao final deste documento}

% Campos opcionais, se houver reinquirição já marcada.
\notificacaoNovoSindicadoDataReinquiricao{00 de mês de 2026}
\notificacaoNovoSindicadoHoraReinquiricao{00h00}
\notificacaoNovoSindicadoLocalReinquiricao{Sala do Sindicante, no quartel do \getOm}

\GerarDocumento{notificacao-novo-sindicado}

% ANEXO K
% ============================================================
% DADOS: ANEXO K - NOTIFICAÇÃO DE DILIGÊNCIAS COMPLEMENTARES
% Arquivo: dados/dados-anexo-k-notificacao-diligencias-complementares.tex
% ============================================================

\novaNotificacaoDiligencias

\notificacaoDiligenciasNumero{001}
\notificacaoDiligenciasOrigem{Sindicante}

% Se ficar vazio, o modelo usa \localData.
% \notificacaoDiligenciasLocalData{Petrolina/PE, 31 de maio de 2026}

\notificacaoDiligenciasDestinatario{\getSindicadoPosto\ \getSindicadoNome\ -- \getSindicadoSU/\getSindicadoOM}

% Se ficar vazio, o modelo usa \autoridadeFuncao.
% \notificacaoDiligenciasAutoridadeFuncao{Comandante do 72º Batalhão de Infantaria de Caatinga}

\notificacaoDiligenciasLocalVista{Sala do Sindicante, no quartel do \getOm}
\notificacaoDiligenciasDescricao{conforme despacho de fls. XX, relativo à realização de diligências complementares}

% Campos opcionais, se houver audiência marcada.
\notificacaoDiligenciasTipoAudiencia{reinquirição do sindicado}
\notificacaoDiligenciasDataAudiencia{00 de mês de 2026}
\notificacaoDiligenciasHoraAudiencia{00h00}
\notificacaoDiligenciasLocalAudiencia{Sala do Sindicante, no quartel do \getOm}

\GerarDocumento{notificacao-diligencias-complementares}

% ANEXO L - segunda aplicação do modelo de juntada
\dataDocumento{8}{maio}{2026}{Salvador/BA}
\chaveCronologica{20260508}
\novaJuntada
\juntadaNumero{2}
\juntadaArquivo{documentos/docjuntada1.pdf}
\juntadaItens{%
  \item Alegações e documentos fictícios apresentados para demonstração do Anexo L.
}
\GerarDocumento{juntada}

% ANEXO M
% ============================================================
% DADOS - ANEXO M - COMPARECIMENTO DE SINDICADO
% Arquivo: dados/dados-anexo-m-comparecimento-sindicado.tex
% ============================================================

% ------------------------------------------------------------
% Data própria do documento
% Deixe em branco para usar a data/local do termo de abertura.
% Formato: \dataDocumento{dia}{mês}{ano}{cidade/UF}
% ------------------------------------------------------------
\dataDocumento{8}{maio}{2026}{CIDADE/UF}
\chaveCronologica{20260508}

\novoComparecimentoSindicado

\comparecimentoSindicadoNumero{001}
\comparecimentoSindicadoOrigem{Sindicante}

% Se não informar, usa \localData.
% \comparecimentoSindicadoLocalData{CIDADE/UF, 31 de maio de 2026}

\comparecimentoSindicadoComandante{Comandante da Subunidade/Seção do Sindicado}
\comparecimentoSindicadoData{05 de junho de 2026}
\comparecimentoSindicadoHora{08h00}
\comparecimentoSindicadoLocal{ORGANIZAÇÃO MILITAR}

\GerarDocumento{comparecimento-sindicado}

% ANEXO N
% ============================================================
% DADOS - ANEXO N - COMPARECIMENTO DE TESTEMUNHA
% Arquivo: dados/dados-anexo-n-comparecimento-testemunha.tex
% ============================================================

\comparecimentoTestemunhaNumero{002}
\comparecimentoTestemunhaOrigem{Sindicante}

% Se não informar, usa \localData.
% \comparecimentoTestemunhaLocalData{Petrolina/PE, 31 de maio de 2026}

% Observação prática:
% - Se a testemunha for militar, endereçar ao comandante.
% - Se for servidor público, endereçar ao chefe.
% - Se for pessoa externa ao Exército, usar modelo de ofício conforme EB10-IG-01.001.
\comparecimentoTestemunhaDestinatario{NOME DA TESTEMUNHA OU COMANDANTE/CHEFE RESPONSÁVEL}
\comparecimentoTestemunhaData{05 de junho de 2026}
\comparecimentoTestemunhaHora{09h00}
\comparecimentoTestemunhaQuartel{72º Batalhão de Infantaria de Caatinga}
\comparecimentoTestemunhaEndereco{Av. Cardoso de Sá, s/n, Petrolina/PE}

\GerarDocumento{comparecimento-testemunha}

% ANEXO O
% ============================================================
% DADOS - ANEXO O - CARTA PRECATÓRIA 1
% Salvar como dados/dados-carta-precatoria1.tex para uso real.
% ============================================================

\novaCartaPrecatoria

\cartaPrecatoriaNumero{001}
\cartaPrecatoriaAutoridadeDeprecada{Comandante do(a) ORGANIZAÇÃO MILITAR}
\cartaPrecatoriaOMDeprecada{ORGANIZAÇÃO MILITAR DEPRECADA}
\cartaPrecatoriaAssunto{inquirição de testemunha}
\cartaPrecatoriaAnexos{cópia da Portaria instauradora; quesitos do sindicante.}
\cartaPrecatoriaFinalidade{inquirição}
\cartaPrecatoriaPessoaOuvida{NOME COMPLETO DA TESTEMUNHA}
\cartaPrecatoriaQualificacaoPessoa{POSTO/GRADUAÇÃO, vinculado ao(à) ORGANIZAÇÃO MILITAR}

\cartaPrecatoriaQuesitos{%
  \item Se tem conhecimento dos fatos objeto da presente sindicância.
  \item Se presenciou diretamente os fatos apurados.
  \item Se tem algo mais a declarar.
}

\GerarDocumento{carta-precatoria}

% ANEXO P
% ============================================================
% DADOS - ANEXO P - TERMO DE INQUIRIÇÃO DE TESTEMUNHA 1
% Salvar como dados/dados-inquiricao-testemunha1.tex para uso real.
% ============================================================

% ------------------------------------------------------------
% Data própria do documento
% Deixe em branco para usar a data/local do termo de abertura.
% Formato: \dataDocumento{dia}{mês}{ano}{cidade/UF}
% ------------------------------------------------------------
\dataDocumento{10}{maio}{2026}{CIDADE/UF}
\chaveCronologica{20260510}

\novaInquiricaoTestemunha

\inquiricaoTestemunhaNumero{1}
\inquiricaoTestemunhaDataAta{quatro dias do mês de maio do ano de dois mil e vinte e seis}

\localInquiricao{Sala do Sindicante, no quartel do \getOm}
\horaInicioInquiricao{09h00}
\horaFimInquiricao{09h30}
\modalidadeInquiricao{presencial}
\inquiricaoHorarioMinutos{540}{570}
\inquiricaoPausaMinutos{0}
\quantidadeTestemunhasDenunciante{2}
\quantidadeTestemunhasSindicado{1}

\testemunhaNome{NOME COMPLETO DA TESTEMUNHA}
\testemunhaPosto{POSTO/GRADUAÇÃO ou qualificação}
\testemunhaNascimento{DATA DE NASCIMENTO}
\testemunhaNaturalidade{NATURALIDADE-UF}
\testemunhaEstadoCivil{ESTADO CIVIL}
\testemunhaFiliacao{NOME DO PAI e NOME DA MÃE}
\testemunhaResidencia{ENDEREÇO COMPLETO}
\testemunhaCelular{(00) 00000-0000}
\testemunhaIdentidade{0000000 ÓRGÃO/UF}
\testemunhaCPF{000.000.000-00}
\testemunhaLocalTrabalho{ORGANIZAÇÃO MILITAR / setor de trabalho}
\testemunhaParentesco{não possui parentesco com o sindicado nem interesse direto no resultado}
\manifestacaoDefesaInquiricao{nada requereu}

\inquiricaoTestemunhaDocumentoBase{os fatos constantes da \getPortaria}
\inquiricaoTestemunhaPerguntaConhecimento{tem conhecimento sobre}
\inquiricaoTestemunhaRespostaConhecimento{sim, tem conhecimento}
\inquiricaoTestemunhaRelato{}
\inquiricaoTestemunhaPerguntas{%
\PergRespAberta{como tomou conhecimento dos fatos}{relatou que...}
\PergResp{presenciou diretamente os fatos}{sim/não, explicando a resposta}
}
\inquiricaoTestemunhaAlgoMais{não}

% Opcional: testemunha da inquirição.
\testemunhaInquiricaoNome{NOME DA TESTEMUNHA DA INQUIRIÇÃO}
\testemunhaInquiricaoPosto{POSTO/GRADUAÇÃO}

\GerarDocumento{termo-inquiricao-testemunha}

% ANEXO Q
% ============================================================
% DADOS - ANEXO Q - TERMO DE INQUIRIÇÃO DE SINDICADO 1
% Salvar como dados/dados-inquiricao-sindicado1.tex para uso real.
% ============================================================

\inquiricaoSindicadoNumero{1}
\inquiricaoSindicadoDataAta{quatro dias do mês de maio do ano de dois mil e vinte e seis}

\localInquiricao{Sala do Sindicante, no quartel do \getOm}
\horaInicioInquiricao{11h30}
\horaFimInquiricao{11h50}

\sindicadoNome{NOME COMPLETO DO SINDICADO}
\sindicadoPosto{POSTO/GRADUAÇÃO}
\sindicadoNascimento{DATA DE NASCIMENTO}
\sindicadoNaturalidade{NATURALIDADE-UF}
\sindicadoEstadoCivil{ESTADO CIVIL}
\sindicadoFiliacao{NOME DO PAI e NOME DA MÃE}
\sindicadoResidencia{ENDEREÇO COMPLETO}
\sindicadoCelular{(00) 00000-0000}
\sindicadoIdentidade{0000000 ÓRGÃO/UF}
\sindicadoCPF{000.000.000-00}

\inquiricaoSindicadoDocumentoBase{\getPortaria}
\inquiricaoSindicadoRespostaConhecimento{sim, tem conhecimento}
\inquiricaoSindicadoRelato{}
\inquiricaoSindicadoPerguntas{%
\PergRespAberta{como o fato aconteceu}{respondeu que...}
\PergResp{teve alguma lesão comprovada}{sim/não, explicando a resposta}
}
\inquiricaoSindicadoAlgoMais{não}
\inquiricaoSindicadoPrazoDefesa{três dias úteis}

% Opcional: testemunha da inquirição.
\testemunhaInquiricaoNome{NOME DA TESTEMUNHA DA INQUIRIÇÃO}
\testemunhaInquiricaoPosto{POSTO/GRADUAÇÃO}

\GerarDocumento{termo-inquiricao-sindicado}

% ANEXO R
% ============================================================
% DADOS: ANEXO R - SUBSTITUIÇÃO DE SINDICANTE
% Arquivo: dados/dados-anexo-r-substituicao-sindicante.tex
% ============================================================

\novaSubstituicaoSindicante

\substituicaoSindicanteNumero{001}
\substituicaoSindicanteOrigem{Sindicante}

% Se ficar vazio, o modelo usa \localData.
% \substituicaoSindicanteLocalData{Petrolina/PE, 31 de maio de 2026}

\substituicaoSindicanteDestinatario{Sr. Comandante do 72º Batalhão de Infantaria de Caatinga}
\substituicaoSindicanteRelatoSucinto{os fatos constantes da portaria instauradora}
\substituicaoSindicanteFls{XX}
\substituicaoSindicanteMotivo{foi verificada situação que recomenda a substituição do sindicante}

\GerarDocumento{substituicao-sindicante}

% ANEXO S
% ============================================================
% DADOS: ANEXO S - TERMO DE ACAREAÇÃO
% Arquivo: dados/dados-anexo-s-acareacao.tex
% ============================================================

% ------------------------------------------------------------
% Data própria do documento
% Deixe em branco para usar a data/local do termo de abertura.
% Formato: \dataDocumento{dia}{mês}{ano}{cidade/UF}
% ------------------------------------------------------------
\dataDocumento{12}{maio}{2026}{CIDADE/UF}
\chaveCronologica{20260512}

\novaAcareacao
\novaAcareacaoDinamica

\acareacaoDataAta{quatro dias do mês de maio do ano de dois mil e vinte e seis}

% Se ficar vazio, o modelo usa: no quartel do(a) \getOm.
% \acareacaoLocal{na Sala do Sindicante, no quartel do \getOm}

\acareacaoPontosDivergencia{descrever objetivamente os pontos divergentes entre os depoimentos}
\adicionarParticipanteAcareacao{NOME COMPLETO DA TESTEMUNHA A}{testemunha}{manteve ou esclareceu sua declaração nos seguintes termos}
\adicionarParticipanteAcareacao{NOME COMPLETO DA TESTEMUNHA B}{testemunha}{manteve ou esclareceu sua declaração nos seguintes termos}
\adicionarParticipanteAcareacao{\getSindicadoNome}{sindicado}{manifestou-se sobre a divergência nos seguintes termos}

% Opcional, se houver advogado presente.
% \advogadoNome{NOME DO ADVOGADO}
% \advogadoOAB{OAB/UF nº 000000}

\dataDocumento{13}{maio}{2026}{Salvador/BA}
\chaveCronologica{20260513}
\GerarDocumento{acareacao}

% ANEXO T
% ============================================================
% DADOS - ANEXO T - TERMO DE ENCERRAMENTO DE INSTRUÇÃO
% Normalmente usa a data principal de dados/dados-sindicancia.tex.
% ============================================================

% Se necessário, redefina a data antes do termo:
% \dia{10}
% \mes{junho}
% \ano{2026}
% \cidade{CIDADE/UF}

% ------------------------------------------------------------
% Data própria do documento
% Deixe em branco para usar a data/local do termo de abertura.
% Formato: \dataDocumento{dia}{mês}{ano}{cidade/UF}
% ------------------------------------------------------------
\dataDocumento{}{}{}{}


\GerarDocumento{termo-encerramento-instrucao}

% ANEXO U
% ============================================================
% DADOS - ANEXO U - VISTA PARA ALEGAÇÕES FINAIS
% Salvar como dados/dados-vista-alegacoes-finais.tex para uso real.
% ============================================================

% ------------------------------------------------------------
% Data própria do documento
% Deixe em branco para usar a data/local do termo de abertura.
% Formato: \dataDocumento{dia}{mês}{ano}{cidade/UF}
% ------------------------------------------------------------
\dataDocumento{15}{maio}{2026}{CIDADE/UF}
\chaveCronologica{20260515}

\novaVistaAlegacoes

\vistaAlegacoesNumero{001}
\vistaAlegacoesPrazo{5 (cinco) dias corridos}
\vistaAlegacoesLocal{Sala do Sindicante, no quartel do \getOm, no período de 08h00 às 12h00}

\GerarDocumento{vista-alegacoes-finais}

% ANEXO V
% ============================================================
% DADOS - ANEXO V - CERTIDÃO
% Arquivo: dados/dados-anexo-v-certidao.tex
%
% Este arquivo define diretamente o texto da certidão.
% Assim, o modelo fica flexível e não exige alteração no sindicancia.cls.
% ============================================================

% ------------------------------------------------------------
% CAMPOS EDITÁVEIS
% ------------------------------------------------------------

\newcommand{\certidaoDataDecursoPrazo}{03 de junho de 2026}
\newcommand{\certidaoDataDecursoPrazoNumerica}{03/06/2026}

% Atenção:
% No modelo enviado havia referência a "29 de junho de 2026", mas isso
% fica cronologicamente incompatível com o decurso de prazo em 03 de junho.
% Ajuste a data abaixo conforme o termo real.
\newcommand{\certidaoDocumentoReferencia}{Termo de Inquirição de Sindicado}
\newcommand{\certidaoDataDocumentoReferencia}{29 de maio de 2026}

\newcommand{\certidaoPrazo}{03 (três) dias úteis}

\newcommand{\certidaoSindicadoNome}{LUCAS VINICIUS DA SILVA RODRIGUES}
\newcommand{\certidaoTestemunhaApresentada}{1º Ten CARVALHO}

\newcommand{\certidaoResultado}{apresentado o \certidaoTestemunhaApresentada\ como testemunha, sem, entretanto, juntar documento extra}

\newcommand{\getCertidaoLocalData}{Petrolina/PE, 03 de junho de 2026}

% ------------------------------------------------------------
% TEXTO DA CERTIDÃO
% ------------------------------------------------------------

\newcommand{\getCertidaoTexto}{%
Certifico que, em \certidaoDataDecursoPrazo, decorreu o prazo concedido por meio do
\certidaoDocumentoReferencia\ de \certidaoDataDocumentoReferencia, concedendo
\textbf{\certidaoPrazo}, a contar da data de inquirição, tendo o sindicado,
\textbf{\certidaoSindicadoNome}, \certidaoResultado.%
}

% ------------------------------------------------------------
% CAIXA DE CIÊNCIA
%
% Para ocultar a caixa, deixe vazio:
% \renewcommand{\getCertidaoCiencia}{}
% ------------------------------------------------------------

\newcommand{\getCertidaoCiencia}{%
Declaro que tenho ciência:

\vspace{0.5em}

Data: \certidaoDataDecursoPrazoNumerica.

\vspace{1em}

Nome:

\vspace{1.2em}

\hrulefill
}

\GerarDocumento{certidao}

% ANEXO W
% ============================================================
% DADOS - ANEXO W - RELATÓRIO
% Arquivo: dados/dados-anexo-w-relatorio.tex
% ============================================================

\relatorioIntroducao{%
A presente sindicância foi instaurada, por determinação do Sr. \getAutoridadePosto\ \getAutoridadeNome, \getAutoridadeFuncao, por intermédio da \getPortaria, para apurar \getObjeto, conforme documentos constantes dos autos, tendo como sindicado \getSindicadoPosto\ \getSindicadoNome.
}

\relatorioDiligencias{%
Com o escopo de reunir elementos probatórios que pudessem esclarecer o fato objeto da presente sindicância, este encarregado houve por bem diligenciar conforme despacho(s) constantes dos autos, tendo sido realizadas as diligências pertinentes, incluindo-se, quando cabível, expedição e recebimento de documentos, inquirições, juntadas e demais atos necessários à elucidação dos fatos.
}

\relatorioParteExpositiva{%
Foi assegurado ao sindicado o direito ao contraditório e à ampla defesa, conforme preconizado nas Instruções Gerais para a Elaboração de Sindicância no Âmbito do Exército Brasileiro -- EB10-IG-09.001.

Da análise de todas as peças que compõem a presente sindicância, restou apurado que: [INSERIR NARRATIVA ORDENADA, COERENTE E CIRCUNSTANCIADA DOS FATOS APURADOS, COM REFERÊNCIA ÀS FOLHAS DOS AUTOS, DEPOIMENTOS, DOCUMENTOS E OUTRAS DILIGÊNCIAS].
}

\relatorioParteConclusiva{%
Em face do exposto, que dos autos consta, e conforme análise realizada na parte expositiva, verifica-se que [INSERIR CONCLUSÃO DO SINDICANTE, CONFORME UMA DAS HIPÓTESES PREVISTAS NO ANEXO W: inexistência de infração, transgressão disciplinar, indícios de crime, dano/extravio de material etc.].

Em consequência, sou de parecer que [INSERIR PROVIDÊNCIA PROPOSTA].
}

% Se não informar, usa \localData.
% \relatorioLocalData{Petrolina/PE, 31 de maio de 2026}

\GerarDocumento{relatorio}

% ANEXO X
% ============================================================
% DADOS - ANEXO X - TERMO DE ENCERRAMENTO
% Arquivo: dados/dados-anexo-x-termo-encerramento.tex
% ============================================================

% Opcional. Se não informar, o modelo usa \getPortaria.
% \termoEncerramentoPortaria{Portaria nº 000, de 00 de mês de 2026, do Sr. Comandante do 72º BI Caat}
\novoTermoEncerramento

\GerarDocumento{termo-encerramento}

% ANEXO Y - encerra o intervalo da foliação do sindicante
% ============================================================
% DADOS - ANEXO Y - DIEx DE REMESSA
% Arquivo: dados/dados-anexo-y-diex-remessa.tex
% ============================================================

% ------------------------------------------------------------
% Data própria do documento
% Deixe em branco para usar a data/local do termo de abertura.
% Formato: \dataDocumento{dia}{mês}{ano}{cidade/UF}
% ------------------------------------------------------------
\dataDocumento{}{}{}{}

\novoDiexRemessa

\diexRemessaNumero{003}
\diexRemessaOrigem{Sindicante}

% Se não informar, usa \localData.
% \diexRemessaLocalData{CIDADE/UF, 31 de maio de 2026}

\diexRemessaDestinatario{\getAutoridadePosto\ \getAutoridadeNome, \getAutoridadeFuncao}
% A quantidade de folhas no assunto do DIEx de remessa é deixada em branco
% no PDF para preenchimento manual no fechamento dos autos.

% Se não informar, usa \getPortaria.
% \diexRemessaReferencia{Portaria nº 000, de 00 de mês de 2026}

% Opcional:
% \diexRemessaTextoComplementar{Informo que os autos seguem em anexo para apreciação e solução por essa autoridade instauradora.}

\GerarDocumento{diex-remessa}

% ANEXO Z - referência da Seção de Apoio Jurídico
% DADOS - ANEXO Z
% Referência da Seção de Apoio Jurídico; não é elaborado pelo sindicante.
\dataDocumento{25}{maio}{2026}{CIDADE/UF}
\chaveCronologica{20260525}
\modoDocumento{padrao}
\solucaoTipo{acolher}
\solucaoParecer{os fatos foram esclarecidos nos termos do relatório}
\solucaoFundamentos{%
  \begin{enumerate}[label=\alph*)]
    \item a instrução observou as formalidades aplicáveis;
    \item as provas foram analisadas de forma comparativa e motivada;
    \item foi assegurado o contraditório quando presente a figura do sindicado.
  \end{enumerate}
}
\solucaoProvidencias{%
  \begin{enumerate}[label=\alph*)]
    \item arquivamento dos autos;
    \item publicação em boletim interno;
    \item ciência aos interessados cabíveis.
  \end{enumerate}
}
\solucaoPublicacaoBI{BI nº 096, de 25 de maio de 2026}
\solucaoNotificados{sindicado e denunciante/ofendido, quando houver}

\GerarDocumento{solucao-sindicancia}

% ANEXO AA - referência da Seção de Apoio Jurídico
\converterRitoParaOrdinario{26 de maio de 2026}{necessidade fictícia de novas provas incompatíveis com o rito sumário}
% DADOS - ANEXO AA
\dataDocumento{12}{maio}{2026}{CIDADE/UF}
\chaveCronologica{20260512}
\modoDocumento{padrao}
\mudancaRitoAutoridade{a autoridade instauradora}
\mudancaRitoFundamento{necessidade de maior profundidade nas investigações}
\mudancaRitoNovasProvas{oitiva de novas testemunhas e realização de perícia técnica}

\dataDocumento{26}{maio}{2026}{Salvador/BA}
\chaveCronologica{20260526}
\GerarDocumento{certidao-modificacao-rito}

\end{document}
